\section{Identificação e Descrição do Problema}
\label{sec:problema}

\subsection{Contexto e Delimitação}
% Descrever o ambiente/contexto onde o problema ocorre.
% Delimitar o escopo: o que está dentro e fora do trabalho.

\subsection{Problema Central}
% Uma pergunta de pesquisa clara ou declaração do problema.
% Ex.: "Como é possível [ação] de forma [característica] para [público]?"

\subsection{Requisitos do Sistema}
% Use duas subseções para requisitos funcionais e não funcionais

\subsubsection{Requisitos Funcionais}
\begin{itemize}
  \item RF01: [Descrição];
  \item RF02: [Descrição];
  \item RF03: [Descrição].
\end{itemize}

\subsubsection{Requisitos Não Funcionais}
\begin{itemize}
  \item RNF01: [Ex.: Desempenho — resposta em menos de 200ms];
  \item RNF02: [Ex.: Usabilidade — interface intuitiva];
  \item RNF03: [Ex.: Portabilidade — executar em Linux e Windows].
\end{itemize}
